\documentclass[12pt,a4paper]{article}

% --- Gói ngôn ngữ & font tiếng Việt ---
\usepackage[utf8]{inputenc}
\usepackage[T5]{fontenc}
\usepackage[vietnamese]{babel}

% --- Căn lề chuẩn văn bản học thuật Việt Nam ---
\usepackage[top=2.0cm, bottom=2.0cm, left=3.0cm, right=2.0cm]{geometry}
\usepackage{setspace}
\onehalfspacing

% --- Gói đồ họa & bảng biểu ---
\usepackage{graphicx}
\usepackage{booktabs}
\usepackage{tabularx}
\usepackage{array}

% --- Định dạng tiêu đề & mục lục ---
\usepackage{titlesec}
\usepackage{tocloft}
\usepackage{hyperref}
\hypersetup{
    colorlinks=true,
    linkcolor=black,
    citecolor=blue,
    urlcolor=blue
}

% --- Tạo khung trang bìa ---
\usepackage{tikz}
\usetikzlibrary{calc}

\begin{document}

% ==================== TRANG BÌA ====================
\begin{titlepage}
\begin{tikzpicture}[remember picture, overlay]
  \draw[line width = 2pt] ($(current page.north west) + (2.0cm,-1.5cm)$)
     rectangle ($(current page.south east) + (-1.5cm,1.5cm)$);
  \draw[line width = 0.5pt] ($(current page.north west) + (2.1cm,-1.6cm)$)
     rectangle ($(current page.south east) + (-1.6cm,1.6cm)$);
\end{tikzpicture}

\begin{center}
    \vspace{0.5cm}
    \textbf{\large BỘ GIÁO DỤC VÀ ĐÀO TẠO}\\[0.2cm]
    \textbf{\large TRƯỜNG ĐẠI HỌC ...}\\[1.5cm]

    \textbf{\Large TIỂU LUẬN MÔN HỌC}\\[0.3cm]
    \textbf{\LARGE KHOA HỌC QUẢN LÝ VÀ QUẢN TRỊ}\\[1.5cm]

    \rule{0.8\textwidth}{0.5pt}\\[0.4cm]
    {\Large \bfseries ĐỀ TÀI:}\\[0.3cm]
    {\LARGE \bfseries PHÂN TÍCH HOẠT ĐỘNG QUẢN LÝ, QUẢN TRỊ TẠI TẬP ĐOÀN FPT VÀ TÁC ĐỘNG CỦA TRÍ TUỆ NHÂN TẠO (AI)}\\[0.4cm]
    \rule{0.8\textwidth}{0.5pt}\\[2.5cm]

    \begin{table}[h]
    \centering
    \begin{tabular}{rll}
        \textbf{Giảng viên hướng dẫn:} & TS./ThS. [Tên giảng viên] & \\
        \textbf{Nhóm sinh viên thực hiện:} & [Họ và tên SV 1] & - MSSV: [MSSV 1] \\
                                           & [Họ và tên SV 2] & - MSSV: [MSSV 2] \\
        \textbf{Lớp học phần:}             & [Mã lớp] & \\
        \textbf{Khóa:}                     & [202... - 202...] &
    \end{tabular}
    \end{table}

    \vfill
    \textbf{Hà Nội / TP.HCM, Tháng .../2026}
\end{center}
\end{titlepage}

% ==================== MỤC LỤC ====================
\newpage
\tableofcontents
\newpage

% ==================== NỘI DUNG CHÍNH ====================

\section*{LỜI MỞ ĐẦU}
\addcontentsline{toc}{section}{LỜI MỞ ĐẦU}
Lý do chọn đề tài, mục tiêu nghiên cứu, đối tượng và phạm vi nghiên cứu, phương pháp thực hiện.

\newpage
\section{TỔNG QUAN VỀ TẬP ĐOÀN FPT}
\subsection{Lịch sử hình thành và phát triển}
Khái quát quá trình từ 1988 đến nay.
\subsection{Lĩnh vực hoạt động cốt lõi và vị thế thị trường}
Khối Công nghệ (FPT Software, FPT IS), Khối Viễn thông (FPT Telecom), Khối Giáo dục (FPT Education).
\subsection{Tầm nhìn, sứ mệnh và giá trị cốt lõi}
Tầm nhìn DC5-135, sứ mệnh kiến tạo hạnh phúc, giá trị văn hóa FPT (Tôn - Đổi - Đồng - Chí - Gương - Sáng).

\newpage
\section{THỰC TRẠNG HOẠT ĐỘNG QUẢN LÝ TẠI FPT}
\subsection{Sơ đồ cơ cấu tổ chức bộ máy quản lý}
Mô hình tổ chức tập đoàn mẹ - con, mô hình quản trị ma trận.
\subsection{Bộ máy điều hành: Vai trò Tổng Giám đốc và các Phó Tổng Giám đốc}
Phân quyền chức năng: Tài chính, Công nghệ, Vận hành, Chuyển đổi số.
\subsection{Cơ chế vận hành và phối hợp giữa Tập đoàn mẹ và các công ty thành viên}
Cơ chế ủy quyền, KPI, OKR trong điều hành vận hành thường nhật.
\subsection{Hạn chế trong công tác quản lý và điểm nghẽn cần cải thiện}
\begin{itemize}
    \item Quy mô nhân sự phình to (>70.000 người), bộ máy nhiều tầng nấc gây giảm tốc độ phản ứng.
    \item Chồng chéo thẩm quyền giữa lãnh đạo sáng lập kỳ cựu và thế hệ lãnh đạo điều hành trẻ.
    \item Độ trễ trao đổi thông tin và xung đột lợi ích cục bộ giữa các công ty con.
\end{itemize}
\subsection{Đề xuất giải pháp hoàn thiện cơ chế quản lý}
Tinh gọn bộ máy, phân cấp triệt để, số hóa toàn diện quy trình tác nghiệp.

\newpage
\section{THỰC TRẠNG HOẠT ĐỘNG QUẢN TRỊ TẠI FPT}
\subsection{Cơ cấu quản trị doanh nghiệp}
Phân tích 3 trụ cột: Đại hội đồng Cổ đông (ĐHĐCĐ), Hội đồng Quản trị (HĐQT), Ban Kiểm soát / Ủy ban Kiểm toán.
\subsection{Vai trò của Hội đồng Quản trị và tính độc lập trong giám sát}
Vai trò Chủ tịch HĐQT, cơ cấu thành viên HĐQT độc lập, các tiểu ban trực thuộc HĐQT.
\subsection{Phân định ranh giới giữa Quản trị (Governance) và Quản lý (Management)}
So sánh thẩm quyền HĐQT (định hướng chiến lược, giám sát, bổ nhiệm) vs Ban Điều hành (thực thi kế hoạch kinh doanh).
\subsection{Đánh giá việc tuân thủ các chuẩn mực quản trị tốt}
Đối chiếu Bộ nguyên tắc Quản trị Công ty Việt Nam (CG Code) và chuẩn mực OECD.

\newpage
\section{TÁC ĐỘNG CỦA TRÍ TUỆ NHÂN TẠO (AI) ĐẾN QUẢN LÝ VÀ QUẢN TRỊ TẠI FPT}
\subsection{Định hướng chiến lược AI của FPT}
Hợp tác NVIDIA xây AI Factory, tích hợp AI vào giải pháp Made by FPT.
\subsection{Tác động của AI đối với hoạt động Quản lý (Management)}
Tự động hóa đánh giá năng suất lập trình, tối ưu quy trình nhân sự, trợ lý ảo điều hành tác nghiệp nội bộ.
\subsection{Tác động của AI đối với hoạt động Quản trị (Corporate Governance)}
Hỗ trợ HĐQT phân tích dữ liệu lớn phục vụ ra quyết định chiến lược, tự động hóa kiểm soát rủi ro và tuân thủ.
\subsection{Thách thức và rủi ro đặt ra từ AI}
Bảo mật dữ liệu, rủi ro thiên vị thuật toán, bài toán đạo đức AI và sự dịch chuyển lực lượng lao động.

\newpage
\section*{KẾT LUẬN VÀ KIẾN NGHỊ}
\addcontentsline{toc}{section}{KẾT LUẬN VÀ KIẾN NGHỊ}
Tóm tắt kết quả nghiên cứu và bài học kinh nghiệm về quản lý, quản trị cho doanh nghiệp Việt Nam.

\newpage
\section*{TÀI LIỆU THAM KHẢO}
\addcontentsline{toc}{section}{TÀI LIỆU THAM KHẢO}
\begin{enumerate}
    \item Báo cáo thường niên Tập đoàn FPT (năm 2023, 2024).
    \item Báo cáo tình hình quản trị công ty Tập đoàn FPT (2023, 2024).
    \item Giáo trình Khoa học Quản lý / Quản trị học.
    \item Bộ nguyên tắc Quản trị Công ty theo thông lệ tốt nhất tại Việt Nam (SSC \& IFC).
\end{enumerate}

\end{document}
